% Triangle Mutiny — Collapse Geometry (v1.2)
% Author: Borealis S. Hedling

\documentclass[11pt]{article}
\usepackage{amsmath, amssymb, amsthm}

\title{Triangle Mutiny: Collapse Geometry (v1.2)}
\author{Borealis S. Hedling}
\date{October 2026}

\begin{document}
\maketitle

\section{Overview}

Collapse geometry formalizes the boundary between stable and unstable
reasoning. A point $p \in M$ lies on the collapse surface $\Sigma$ when
any structural invariant fails: curvature continuity, spectral continuity,
modal adjacency, or drift compatibility.

This module depends on the primitives defined in \texttt{core-math.tex}.

\section{Collapse Surface}

The collapse surface is defined as:
\[
\Sigma
  = \{ p \in M :
        C_K(p) \lor
        C_{\mathrm{spec}}(p) \lor
        C_{\mathrm{modal}}(p) \lor
        C_{\mathrm{drift}}(p)
     \}.
\]

Each predicate corresponds to a structural violation.

\section{Curvature Collapse}

Curvature collapse occurs when curvature continuity fails:
\[
C_K(p)
  := \exists \epsilon > 0 \text{ such that }
     \lim_{q \to p} |K(q) - K(p)| > \epsilon.
\]

\subsection{Lemma: Curvature Discontinuity Implies Collapse}
If $C_K(p)$ holds, then $p \in \Sigma$.

\begin{proof}
Immediate from definition of $\Sigma$.
\end{proof}

\section{Spectral Collapse}

Spectral collapse occurs when spectral continuity fails:
\[
C_{\mathrm{spec}}(p)
  := \exists \epsilon > 0 \text{ such that }
     \lim_{q \to p}
       \|\mathrm{Spec}(L(q)) - \mathrm{Spec}(L(p))\|
       > \epsilon.
\]

\subsection{Lemma: Spectral Discontinuity Implies Collapse}
If $C_{\mathrm{spec}}(p)$ holds, then $p \in \Sigma$.

\begin{proof}
Direct from definition.
\end{proof}

\section{Modal Collapse}

Modal collapse occurs when a transition violates adjacency:
\[
C_{\mathrm{modal}}(p)
  := (\mu(p^-), \mu(p^+)) \notin E.
\]

\subsection{Lemma: Invalid Modal Transition Implies Collapse}
If $C_{\mathrm{modal}}(p)$ holds, then $p \in \Sigma$.

\begin{proof}
By definition of modal adjacency constraints.
\end{proof}

\section{Drift Collapse}

Drift collapse occurs when drift compatibility fails:
\[
C_{\mathrm{drift}}(p)
  := [X,[Y,Z]] + [Y,[Z,X]] + [Z,[X,Y]] \neq 0.
\]

This corresponds to Jacobi identity violation.

\subsection{Lemma: Jacobi Failure Implies Collapse}
If $C_{\mathrm{drift}}(p)$ holds, then $p \in \Sigma$.

\begin{proof}
Jacobi identity is required for operator stability.
\end{proof}

\section{Composite Collapse Conditions}

A point collapses if any structural invariant fails:
\[
p \in \Sigma
  \iff
  C_K(p)
  \lor C_{\mathrm{spec}}(p)
  \lor C_{\mathrm{modal}}(p)
  \lor C_{\mathrm{drift}}(p).
\]

\section{Collapse Neighborhood}

Define the collapse neighborhood:
\[
U(\Sigma)
  = \{ q \in M : \mathrm{dist}(q, \Sigma) < \delta \}.
\]

Points in $U(\Sigma)$ are considered near-collapse and require spectral
and drift monitoring.

\section{Collapse Severity}

Define collapse severity:
\[
S_{\Sigma}(p)
  = \alpha |K(p^+) - K(p^-)| 
  + \beta \delta_{\mathrm{spec}}(p)
  + \gamma \Delta_{\mathrm{modal}}(p)
  + \eta \|J(p)\|,
\]
where $J(p)$ is the Jacobi failure magnitude.

Severity is used in stability envelope classification.

\end{document}
